\section{Conclusion}

APT29, a Russian based and foreign intelligence associated hacker group focused on digital stealth infiltration and cyber-espionage. Characterised by abnormally skilled operators with global influence. Unlike other hacker groups, they've shown to act more out of state interest than monetary gain. Operations such as the HAMMERTOSS backdoor show their effectiveness at understanding human behaviour and integrating it into their attacks and infiltrations. This means they make an especially dangerous threat utilising tools like social engineering and espionage to tailor their attacks towards individuals or specific high profile institutions.

Protecting against threats like these requires holistic approaches rather than specific low level measures. Companies need to consider their internal cyber security culture, and verification pipelines. Making sure that employees are able to detect potentially novel and targeted forms of attacks, and/or process data through the correct internal bodies who make sure it's safe. Due to the scale of the threat it is also important for companies to be able to detect and deal with it once they *have* been infiltrated. Managers need to create a culture in which reporting incidents is the norm rather than having employees avoid it in fear of repercussions. They also need to perform routine internal evaluation to make sure they are keeping up to date with policies and its effectiveness.